\documentclass{iutmscthesis}

% The iutbscthesis uses biblatex for bibliography management.
% You can learn the basics of biblatex from:
% https://www.overleaf.com/learn/latex/Bibliography_management_with_biblatex
\addbibresource{citations.bib}

% The complete title of your thesis, mandatory
\title{Lightweight Trap-Based Free-Rider Detection and Mitigation in Federated Learning}

% Name and Student ID of the author, mandatory
\author{Aashnan Rahman}{231041006}

% Use the \supervisor command to define the supervisor with name, designation, and department
% There must be one and only one supervisor, mandatory
\supervisor{Dr.\ Md.\ Azam Hossain}{Associate Professor}{Department of Computer Science and Engineering}

% Use the \cosupervisor command to add a cosupervisor name, designation, 
% department, and organization/institution. The cosupervisor can be from
% an outside university or organization. There can be at max one cosupervisor,
% optional. If there is no cosupervisor, just remove the \cosupervisor command.
\cosupervisor{Dr.\ Mike Hunt}{Professor}{Department of Electrical and Computer Engineering}{Imaginary University of Technology, Texas, US}

% Name, designation, and department of the head, as an ex-officio member, mandatory
\depthead{Dr.\ Justin Case}{Professor}{Department of Computer Science and Engineering}

% Name, designation, and department of the first internal member, mandatory
\member{Dr.\ Benjamin Dover}{Assistant Professor}{Department of Computer Science and Engineering}

% Name, designation, and department of the second internal member, optional
% If there is no second internal member, just remove the \membertwo command
\membertwo{Mr. Noah Pinion}{Assistant Professor}{Department of Computer Science and Engineering}

% Name, designation, department, and organization/institution of the external member, mandatory
\external{Toby Lerone}{Professor}{Department of Computer Science and Engineering}{Bjork University, UK}

% Specific degree in masters, default is `Master of Science in Computer Science and Engineering`
\degree{Master of Science in Computer Science and Engineering}

% Complete department name, mandatory
\department{Department of Computer Science and Engineering}

% Academic year in range separated by endash, mandatory
\academicyear{2024--25}

% Date (day, month name, and year) when the presentation took place, mandatory
\defensedate{04}{July}{2026}

\begin{document}

%-------------------------- FRONT MATTER --------------------------%

% Generate cover page of the report to be printed on the hard cover
\coverpage

% Switch to roman page numbering
\pagenumbering{roman}

% Generate the title leaf, mandatory
\titlepage

% Auto generate declaration of candidate based on previous info, mandatory
\certificateofapproval
\declarationofcandidate

% Dedicate report to someone. You may add more sentences to it.
%\dedicatedto{my parents, my wife, and my daughter without whom this degree would have been completed two years earlier}

\dedicatedto{To be added later}
\tableofcontents
\listoffigures
\listoftables

\clearpage

% Add abbreviations within this environment using the \abbr command
\begin{abbreviations}
    \abbr{FL}{Federated Learning}
    \abbr{FR}{Free Rider}
    \abbr{IID}{Independent and Identically Distributed}
    \abbr{MAD}{Median Absolute Deviation}
    \abbr{FedAvg}{Federated Averaging}
\end{abbreviations}

% Any acknowledgement that you want to write goes in the following environment
\begin{acknowledgement}

I am profoundly grateful to my supervisor, Dr.\ Lou Natic, for his endless patience, insightful critiques, and for always knowing just when to say, ``Have you tried turning it off and on again?'' His expertise and encouragement were instrumental in the completion of this thesis.

A special thanks to my co-supervisor, Dr.\ Sue Permann, who provided invaluable support and the occasional pep talk disguised as a lecture on time management. Her unwavering belief in my work kept me going, even when the only thing I believed in was the power of procrastination.

Finally, to my parents, whose love and support have been a constant source of strength. Thank you for always being there, for pretending to understand my research when I rambled on, and for not asking too many questions about why I was still not employed after all these years.

To all of you, I extend my heartfelt appreciation and a promise: next time, I will take fewer detours through the land of streaming services.
\end{acknowledgement}

% Write your abstract here.
\begin{abstract}
Federated learning allows clients to benefit from a shared model without centralizing raw data, but the server cannot directly verify whether a submitted update resulted from genuine local training. This thesis presents a lightweight server-side mechanism that secretly replaces the global model with a randomized trap model for selected clients and evaluates the response using robust update-norm and layer-profile statistics. The proposed method combines a two-trap coverage sweep, candidate confirmation, frequent re-probing of suspicious clients, rotating anchor panels, continued surveillance, cumulative penalties, and evidence-based rehabilitation. Detection uses only returned model parameters; client-reported accuracy, loss, training time, and resource use are not trusted inputs.

The method was evaluated on MNIST and CIFAR-10 with 100 clients and 100 communication rounds under IID and Dirichlet non-IID partitions ($\alpha=0.5$), four free-rider strategies, attacker shares of 30\% and 40\%, and repeated random seeds. The best-observed summaries achieve 100\% recall in every displayed configuration. Precision is 93.75--100\% on MNIST and 98.36--100\% on IID CIFAR-10. The new CIFAR-10 30\% batch removes every free rider in all sixteen runs, but its non-IID executions also remove 7--18 honest clients each; the best-observed non-IID 30\% precision is therefore only 71.43\%. Earlier MNIST repetitions additionally expose an adaptive FR4 contamination limit. Final global accuracy remains 98.82--98.95\% on MNIST and approximately 74.48--77.89\% on CIFAR-10.

The findings show that secret statistical probing is effective and inexpensive for replay, random-update, and historical-average attacks, but that an unflagged client is not necessarily a trustworthy reference and robust statistics alone do not guarantee low false-positive rates under heterogeneous data. The method is therefore reported as a strong lightweight baseline with a clearly identified non-IID limitation rather than as a universal proof of client computation.
\end{abstract}

%-------------------------- MAIN BODY --------------------------%
% Switch to arabic page numbering
\pagenumbering{arabic}

\chapter{Introduction}


\section{Introduction to Federated Learning and the Free Rider Problem}
\label{sec:intro-info}

Federated Learning (FL) has emerged as a transformative paradigm for distributed machine learning that enables multiple clients such as mobile devices, edge servers, or organizational silos to collaboratively train a shared global model without exposing their raw, locally stored data \cite{mcmahan2017communication,kairouz2021advances}. Unlike traditional centralized machine learning, where data must be aggregated on a single server, FL operates on a fundamentally different principle: \emph{the model travels to the data, not the other way around}.

In a typical FL architecture, a central server initializes a global model and distributes it to a set of participating clients. Each client independently trains the model on its local dataset for a fixed number of epochs and computes a local model update (e.g., gradients or updated weights). These updates are transmitted back to the server, which aggregates them—commonly via the Federated Averaging (FedAvg) algorithm—to produce an improved global model \cite{mcmahan2017communication}. This process is repeated over multiple communication rounds until the global model converges to a satisfactory performance level.

This architecture offers several compelling advantages:

\begin{itemize}
    \item \textbf{Privacy Preservation:} Since raw data never leaves the client device, FL inherently reduces the risk of data leakage and aligns with data protection regulations such as GDPR and HIPAA.
    \item \textbf{Reduced Communication Overhead:} Only model parameters or gradients are exchanged, rather than entire datasets, which is particularly advantageous in bandwidth-constrained environments.
    \item \textbf{Decentralized Data Utilization:} FL enables organizations and individuals to benefit from collaborative model training without relinquishing control over sensitive data.
    \item \textbf{Scalability:} FL systems can, in principle, incorporate thousands to millions of participating clients, as seen in large-scale mobile deployments such as Google's keyboard prediction systems \cite{hard2018federated,bonawitz2019towards}.
\end{itemize}

Privacy is not merely a desirable feature of FL it is often the primary motivation for its adoption. In domains such as healthcare, finance, and personal mobile applications, the ability to train robust models while keeping sensitive data on-device is a critical enabler of machine learning adoption in regulated and privacy-sensitive contexts.

However, the very decentralization and privacy guarantees that make FL attractive also introduce a unique vulnerability: the server has limited visibility into what each client is actually doing locally. This opacity creates an opportunity for malicious or self-interested clients to exploit the system through what is known as \textbf{Free Rider (FR) behavior}. A free rider is a participant that contributes little or no genuine computational effort toward local model training, yet still submits an update to the server in order to receive the benefits of the collaboratively trained global model—effectively benefiting from the effort of honest participants while incurring minimal cost. Because free riders can fabricate or manipulate updates that superficially resemble legitimate contributions, they are often difficult to distinguish from genuine, well-behaved clients using standard aggregation mechanisms alone.


%\section{Introductory Information}



\section{Scope and Motivation}
\label{sec:scope-motivation}

To understand why free rider behavior poses a serious threat to FL systems, consider a simple real-world analogy. Suppose a group project involves ten students who are collectively assigned a task, and their final grade depends on the group's combined output. If six students genuinely contribute their effort while the remaining four merely submit superficial or copied work without doing any real analysis, the final result may still appear complete and acceptable to an external evaluator who only inspects the submitted output. Yet the four free-riding students receive the same reward (the grade) as the diligent six, despite having contributed little of value. Over time, this dynamic discourages honest participants, degrades the overall quality of the group's output, and undermines the fairness of the collaborative process.

This scenario maps directly onto FL systems. Free-riding clients can submit updates that are statistically or structurally similar to genuine updates—for instance, by adding small perturbations to the previous global model, reusing outdated updates, or applying minimal or synthetic training—without performing meaningful local training on their own data. Because the FL server typically aggregates updates based on volume or simple heuristics (e.g., number of local samples reported by the client) rather than verifying the training process itself, such free riders can successfully "blend in" with genuine contributors. They benefit from the improved global model at every round while incurring negligible computational, energy, or data costs.

The consequences of unmitigated free-rider behavior in FL systems are significant:

\begin{itemize}
    \item \textbf{Degraded Model Quality:} A high proportion of free riders can slow convergence or degrade the accuracy of the global model, especially in non-IID data settings where genuine diverse contributions are essential.
    \item \textbf{Unfair Resource Allocation:} In incentive-based FL systems, free riders may unfairly claim rewards (monetary or otherwise) intended to compensate genuine contributors, discouraging honest participation over time.
    \item \textbf{Security and Trust Erosion:} Persistent free riding undermines trust in the FL ecosystem, potentially disincentivizing high-quality clients from participating altogether.
    \item \textbf{Scalability of Malicious Behavior:} As FL systems scale to include thousands of clients, manual or centralized verification of each client's genuine contribution becomes computationally and practically infeasible, making the free rider problem increasingly urgent.
\end{itemize}

Given the growing deployment of FL in real-world, large-scale, and often adversarial environments, addressing free rider behavior is not merely an academic exercise but a practical necessity for ensuring the long-term viability, fairness, and robustness of federated systems.


\section{Problem Statement}
\label{sec:problem-statement}

The core problem addressed in this thesis can be formally framed as follows:

\begin{quote}
\textit{In a Federated Learning system involving $N$ clients participating over $T$ communication rounds, a subset of clients $\mathcal{F} \subset \mathcal{N}$ may behave as free riders—submitting model updates that do not reflect genuine training on locally available data, while still appearing statistically and structurally consistent with updates submitted by honest clients. The central server, having no direct access to client-side data or training logs, must identify or mitigate the influence of these free-riding clients using only the information available through the exchanged model updates, aggregation metadata, and system-level observations, without compromising the privacy guarantees that FL is designed to provide.}
\end{quote}

This problem is inherently challenging because it sits at the intersection of two competing goals: (1) verifying the authenticity and quality of client contributions, and (2) preserving the privacy-preserving nature of the FL paradigm, which explicitly avoids direct inspection of client data or computation. A free rider detection or mitigation mechanism that requires access to raw client data, imposes excessive computational burden, or fails to generalize across dynamic and heterogeneous client populations is, in practice, unsuitable for real-world FL deployments.

\section{Research Challenges}
\label{sec:research-challenges}

Detecting and mitigating free rider behavior in FL systems presents several interrelated challenges:

\begin{enumerate}
    \item \textbf{Indistinguishability of Updates:} Free riders can craft updates through techniques such as noise injection, weight perturbation, or reuse of previous global models—that closely mimic the statistical properties of genuine updates, making detection based on superficial inspection unreliable.

    \item \textbf{Absence of Ground Truth:} The server cannot directly observe a client's local training process, dataset, or computational effort, making it difficult to establish a reliable baseline against which suspicious updates can be compared.

    \item \textbf{Non-IID Data Distributions:} In realistic FL settings, client data is often non-independent and non-identically distributed (non-IID), a well-known source of optimization instability and client-update diversity in FL \cite{zhao2018federated,kairouz2021advances}. Legitimate updates from clients with unusual or minority data distributions may appear anomalous, increasing the risk of falsely flagging honest clients as free riders.

    \item \textbf{Privacy Constraints:} Any detection mechanism must operate without requiring access to raw client data or intrusive monitoring, as doing so would violate the fundamental privacy guarantees that motivate the use of FL in the first place. This constraint becomes even stronger in deployments that rely on secure aggregation, where the server is deliberately prevented from inspecting individual client updates \cite{bonawitz2017practical}.

    \item \textbf{Scalability:} FL systems may involve thousands to millions of clients. Detection mechanisms that rely on computationally expensive verification, pairwise comparisons, or centralized auditing do not scale efficiently to such large populations.

    \item \textbf{Dynamic and Adaptive Adversaries:} Free riders may adapt their strategies over time to evade detection (e.g., alternating between genuine and free-riding behavior), requiring mitigation techniques that are robust to adaptive and evolving attack patterns rather than static, one-time detection.

    \item \textbf{Trade-off Between Robustness and Model Utility:} Overly aggressive filtering or penalization of suspicious clients risks excluding legitimate contributions, particularly from clients with limited resources or atypical data, thereby harming overall model utility and fairness.
\end{enumerate}

Several existing approaches have attempted to address aspects of this problem, but each carries notable limitations:

\begin{itemize}
    \item \textbf{Hardware-based methods}, such as trusted execution environments (TEEs) or secure hardware attestation, can verify that computation genuinely occurred on a client device. However, these methods require specialized hardware support, are costly to deploy at scale, and are often infeasible for heterogeneous, resource-constrained edge devices.

    \item \textbf{Proof-of-work approaches}, inspired by blockchain consensus mechanisms, require clients to demonstrate computational effort before their updates are accepted. While effective in principle, these methods impose significant computational and energy overhead, which is particularly problematic for battery-constrained mobile or IoT devices, and does not directly guarantee that the effort corresponds to genuine, useful model training.

    \item \textbf{Server-side verification} techniques attempt to validate updates by re-executing or partially replicating client-side computation. These approaches, however, often require access to client data or proxies thereof, directly conflicting with FL's privacy-preserving objectives and with secure-aggregation-style deployments \cite{bonawitz2017practical}, and introduce substantial computational burden on the server.

    \item \textbf{Computation checking} methods analyze metadata such as training time, number of local epochs, or resource usage logs reported by clients. Such self-reported metrics are inherently susceptible to falsification by malicious clients, limiting their reliability as a standalone defense.

\end{itemize}

Collectively, these limitations—high cost, privacy compromise, poor scalability, and susceptibility to adaptive adversaries—highlight the need for a fundamentally different approach: one that is lightweight, privacy-preserving, scalable to large and dynamic client populations, and robust against evolving free-riding strategies.


\section{Contributions}
\label{sec:contributions}

Motivated by the limitations of existing free rider mitigation strategies, this thesis proposes a lightweight, privacy-preserving, and scalable framework for detecting and mitigating free rider behavior in Federated Learning systems. The key contributions of this work are summarized as follows:

\begin{enumerate}
    \item \textbf{A Privacy-Preserving Detection Mechanism:} We design a detection approach that relies solely on information already available at the server through the standard FL communication protocol—such as update characteristics and aggregation-level statistics—without requiring access to raw client data, specialized hardware, or intrusive auditing procedures.

    \item \textbf{Lightweight Computational Overhead:} Unlike proof-of-work or server-side re-execution methods, the proposed approach introduces minimal additional computational and communication cost, making it practical for deployment in resource-constrained and large-scale FL environments.

    \item \textbf{Robustness to Dynamic and Adaptive Free Riders:} The proposed method is designed to adapt to evolving free-riding strategies over multiple communication rounds, rather than relying on static, one-time verification, thereby improving robustness against adaptive adversarial behavior.

    \item \textbf{Preservation of Model Utility for Honest Clients:} By carefully balancing detection sensitivity against the risk of false positives, particularly in non-IID settings, the proposed approach aims to mitigate free rider influence without unfairly penalizing legitimate clients with atypical but genuine data distributions.

    \item \textbf{Scalability to Large Client Populations:} The proposed framework is designed with computational and communication efficiency in mind, enabling its application to FL systems involving large and dynamically changing sets of participating clients.

    \item \textbf{Empirical Validation and Failure Analysis:} We evaluate four free-rider strategies on MNIST under IID and Dirichlet non-IID partitions, two attacker shares, and two random seeds. Beyond precision, recall, F1, model accuracy, and runtime, the analysis traces false removals to individual clients and identifies a reproducible FR4 failure caused by statistical and anchor-set contamination at 40\% non-IID attackers. No unimplemented comparison with external baselines is claimed.
\end{enumerate}

Collectively, these contributions aim to advance the practical deployability of Federated Learning systems by addressing one of its most persistent and underexplored vulnerabilities, offering a solution that preserves the core privacy and efficiency benefits that make FL attractive in the first place.


\subsection{Organization}
Chapter~1 introduces the free-rider problem, research challenges, and contributions. Chapter~2 reviews statistical, learned, audit-based, privacy-attack-based, and hardware-assisted defenses. Chapter~3 specifies the proposed method, including its threat model, two-signal detector, probing schedule, client states, penalties, rehabilitation rule, and aggregation policy. Chapter~4 presents the experimental protocol and analyzes IID and non-IID results, model utility, runtime, false positives, and the FR4 contamination failure. Chapter~5 summarizes the findings, limitations, and priorities for future work.

\chapter{Related Work}
\label{sec:related-work}

\section{Background}
Research on free-rider (FR) behavior in Federated Learning (FL) can be organized into five broad categories: (1) foundational work that formalizes free-rider attacks and analyzes their effect on convergence, (2) increasingly stealthy attack constructions that exploit the geometric and statistical properties of honest updates, (3) server-side detection and audit-based mitigation mechanisms, (4) detection methods that repurpose privacy-attack techniques, and (5) hardware-based and cryptographic verification schemes. This section reviews each category in turn, compares the underlying assumptions and effectiveness of representative methods, and concludes by identifying the gap that motivates the approach proposed in this thesis.

\section{Foundational Characterization of Free-Rider Attacks}
\label{subsec:foundational}

Fraboni et al. \cite{fraboni2021freerider} provided the first theoretical and experimental analysis of free-rider attacks against iterative parameter-aggregation schemes such as FedAvg and FedProx. Modeling stochastic gradient descent as a continuous-time process, the authors formally proved that a free rider's fabricated update can be constructed to provably converge toward the aggregated model produced by fair participants, meaning the attacker obtains a fully-trained global model without contributing any data. They first showed that the most naive attack—simply returning the previous global model unchanged—is trivially detectable by basic server-side checks (e.g., zero update norm). To evade such checks, they proposed a \emph{disguised free-rider} strategy that adds time-decaying Gaussian noise to the previous global model, with the noise variance calibrated to the empirical standard deviation of the global gradient so that the fabricated update statistically mimics the natural shrinkage of SGD updates as training converges. Their experiments across both IID and non-IID settings showed that this disguised attack is substantially harder to detect than the naive variant, and they concluded with practical recommendations for FL deployments in sensitive, high-value domains such as healthcare.

Complementing this, Lin et al. \cite{lin2019freeriders} were the first to formally introduce the notion of free-rider attacks in FL and to jointly study attack construction and detection. They proposed the \emph{delta weights} attack, in which a free rider constructs its fake update by adding calibrated noise to the previous round's global gradient rather than to the model weights directly, producing updates that more closely resemble genuine gradient-based contributions than naive random-weight submissions. On the defense side, they proposed \textbf{STD-DAGMM}, a detection method that augments a Deep Autoencoding Gaussian Mixture Model (DAGMM) with the standard deviation of each client's update as an additional feature, framing free-rider detection as a high-dimensional anomaly-detection problem. They further extended their analysis to scenarios involving multiple simultaneous free riders and the presence of differential privacy noise, showing that the interaction between privacy-preserving noise and free-rider disguise strategies further complicates detection—an open problem they explicitly flagged for future research.

Together, these two works established that (i) free-riding is a concretely realizable and theoretically grounded threat rather than a hypothetical one, and (ii) naive server-side aggregation, and even simple statistical checks, are insufficient once an attacker calibrates its fabricated update to mimic the stochastic behavior of genuine local training.

\section{Advanced and Stealthy Free-Rider Attacks}
\label{subsec:advanced-attacks}

Building directly on the disguising strategies above, Zhu et al. \cite{zhu2023advanced} introduced a more sophisticated \emph{advanced free-rider attack} that explicitly targets the geometric properties exploited by feature-based defenses. The authors analyzed the expected cosine similarity and $\ell_2$-norm relationships between honest client updates and the converging global model in IID settings, deriving closed-form approximations for how these quantities evolve as training progresses. Using these derived relationships, their attack injects calibrated Gaussian noise into only a \emph{subset} of the global gradient's parameters, with the noise magnitude explicitly scaled so that the resulting fake update reproduces the cosine-similarity and norm statistics expected of a genuine client update at that point in training. Because this attack is constructed to match the exact geometric signature that earlier detection mechanisms (including cosine-similarity and $\ell_2$/standard-deviation-based checks) rely on, it was shown—both in the original work and in later independent evaluations \cite{recasens2024frida}—to substantially degrade the effectiveness of feature-based detectors, particularly in IID scenarios where honest client updates are naturally similar to one another and thus offer the attacker less statistical cover to hide \emph{behind} but also less distinctiveness to be caught \emph{by}.

This line of work demonstrates a critical limitation of static, feature-based defenses: any detection rule based on a fixed set of update statistics (norm, standard deviation, cosine similarity) can, in principle, be reverse-engineered and matched by an adaptive attacker with knowledge of the aggregation dynamics. It motivates the need for detection mechanisms that target properties of genuine training which are fundamentally difficult to fabricate, rather than statistical proxies that a sufficiently informed adversary can approximate.

\section{Server-Side Detection and Audit-Based Mitigation}
\label{subsec:server-side}

A second line of work designs server-side detection pipelines that combine anomaly detection with auxiliary contribution or reputation signals. Wang et al. proposed \textbf{FRAD} \cite{wang2024frad}, a free-rider attack detection mechanism targeted at FL deployments in the Artificial Intelligence of Things (AIoT), where device heterogeneity is pronounced. FRAD combines a DAGMM-based anomaly detector (extending the STD-DAGMM approach of Lin et al.) with an explicit \emph{contribution model} that estimates each device's expected contribution from its computing resources, communication cost, and data quality, and a \emph{reputation model} based on the PageRank algorithm that iteratively re-weights clients according to their historical behavior. By fusing statistical anomaly detection with resource- and reputation-aware contribution modeling, FRAD is explicitly designed to select benign nodes for participation even under conditions of information asymmetry between the server and heterogeneous IoT clients. However, because its detection backbone remains a DAGMM-style autoencoder, it inherits the same scalability limitation identified in later work \cite{recasens2024frida}: the memory and computational cost of training an autoencoder over every client's full (flattened) update grows with both model size and the number of clients, which becomes prohibitive for larger neural networks or large client populations.

Wang et al. also proposed \textbf{PASS} \cite{wang2023pass}, a Parameter Audit-based Secure and fair FL Scheme, motivated by the observation that prior defenses generally assumed adversaries constitute a minority of clients and focused narrowly on \emph{anonymous} free riders (clients lacking data or compute) while losing effectiveness against \emph{selfish} free riders (clients that possess adequate resources but strategically withhold them). PASS combines a privacy-preserving auditing strategy—injecting weak differential-privacy noise via a Gaussian mechanism together with a parameter-pruning step to protect client data during the audit—with a novel contribution-evaluation method used to assess each client's genuine performance. The scheme's experimental evaluation demonstrated that PASS achieves privacy protection comparable to state-of-the-art baselines with limited accuracy loss, attains a higher defense success rate, lower false-positive rate, and higher F1-score than prior methods against both anonymous and selfish free-rider attacks, and—notably—remains effective even when adversarial clients constitute a \emph{majority} of the population (F1-scores of 89\% and 87\% against anonymous and selfish attacks, respectively, even when adversaries exceed 50\% of clients). This robustness to majority-adversary settings distinguishes PASS from earlier defenses, though its parameter-auditing procedure introduces non-trivial computational overhead at the server compared to lightweight statistical baselines.

Collectively, FRAD and PASS demonstrate that meaningful detection is achievable using only server-observable signals augmented with lightweight contribution modeling, aligning with FL's privacy-preserving ethos. However, both rely on carefully engineered statistical or audit pipelines whose effectiveness depends on assumptions (e.g., DAGMM's density-estimation quality, or the calibration of PASS's audit noise) that may require retuning as data distributions, model architectures, and free-rider strategies evolve.

\section{Privacy-Attack-Based Detection}
\label{subsec:privacy-attack-based}

A more recent line of work, exemplified by Recasens et al.'s \textbf{FRIDA} (Free-Rider detection using privacy Attacks) \cite{recasens2024frida}, departs from indirect statistical proxies and instead repurposes \emph{privacy-attack} methodologies to directly verify whether a client's update reflects genuine training. The key insight is that prior detection families—feature-based statistics (norm, standard deviation, cosine similarity), DAGMM-based anomaly detection (including FRAD), log-based auditing, and incentive/reputation mechanisms—all rely on \emph{indirect effects} of free-riding rather than on the root cause: the absence of genuine training on real client data. FRIDA instead proposes four detection mechanisms built on two families of privacy attack. The first family adapts \emph{membership inference attacks} (MIA), originally studied as a way to infer whether a record was part of a model's training data \cite{shokri2017membership}: the server distributes a small shared ``canary'' dataset that honest clients are required to additionally train on each round; a \emph{loss-based} detector then flags clients whose canary-set loss fails to decrease as expected, while a \emph{cosine-based} detector compares the cosine similarity of a client's gradient with canary-sample gradients against non-canary samples via a statistical test, flagging clients for whom the two distributions are indistinguishable (indicating no canary training occurred). The second family adapts \emph{property inference attacks} (PIA), which exploit information leakage from collaborative model updates to infer aggregate properties of a participant's local data \cite{melis2019exploiting}, to estimate each client's local label distribution from its submitted gradients without requiring any shared dataset or additional client-side training; a \emph{consistency-based} detector flags clients whose inferred label distribution changes anomalously across rounds, while a \emph{diversity-based} detector flags clients whose inferred distribution is best explained as a combination of pure noise and the global label distribution—the signature of a fabricated, rather than genuinely trained, update.

The authors' extensive evaluation across CIFAR-10 and CIFAR-100 with AlexNet and VGG-19 architectures, under both IID and non-IID data (Dirichlet-distributed with $\alpha \in \{0.1, 1.0\}$), and against all three attack strategies discussed above (Fraboni et al., Lin et al., and Zhu et al.), showed that FRIDA matches or outperforms feature-based and DAGMM-based baselines in nearly every setting, with especially pronounced advantages in IID scenarios and on more complex datasets—precisely the regimes in which the advanced attack of Zhu et al. is most effective at evading feature-based detectors. Notably, PIA-based detection was shown to remain robust even when honest clients apply local differential privacy to their updates, since it exploits dataset-level statistical properties rather than fine-grained individual-level signals, whereas MIA-based detection degrades sharply under such privacy noise. However, FRIDA's design assumes the server has visibility into individual client updates, which is fundamentally in tension with secure-aggregation-based privacy protection; the authors propose a two-server architecture as a partial mitigation but acknowledge this adds system complexity, and the cosine-based MIA variant was shown to be evadable by an adaptive, resource-endowed free rider willing to train specifically on the canary set. The MIA-based variants also introduce non-trivial per-round computational cost relative to lightweight feature-based statistics, and PIA-based detection shifts additional computational burden onto the server.

\section{Hardware-Based and Cryptographic Solutions}
\label{subsec:hardware-based}

In parallel with statistical and inference-based defenses, several works have explored hardware-based and cryptographic approaches. Trusted Execution Environments (TEEs), such as Intel SGX, have been proposed to provide attestable, tamper-evident proof that local training genuinely executed on a client device. Similarly, blockchain-based FL frameworks have incorporated proof-of-work or proof-of-training consensus mechanisms that require clients to demonstrate verifiable computational effort before their updates are accepted into aggregation, and secure-aggregation or zero-knowledge-proof-based schemes have been proposed to protect or verify client-side computation without exposing raw client data or full individual updates \cite{bonawitz2017practical}.

While such approaches can offer strong, cryptographically verifiable guarantees against fabricated updates, they carry substantial practical costs. TEE-based solutions require specialized hardware support that is unavailable across much of the heterogeneous device landscape typical of real-world FL (older mobile devices, low-cost IoT sensors), limiting deployability. Proof-of-work mechanisms impose significant computational and energy overhead that is particularly problematic for battery-constrained edge devices and works directly against FL's goal of enabling participation from resource-limited clients. Zero-knowledge-proof schemes, while avoiding the hardware-dependency problem, typically introduce substantial proof-generation and verification overhead that scales poorly with model size. None of these approaches inherently addresses the non-IID, dynamic nature of real-world client populations, and all introduce deployment complexity that has, to date, limited their adoption outside specialized or high-value settings.

\section{Other Notable Approaches}
\label{subsec:other-approaches}

Beyond the categories above, several related lines of work address free riding indirectly through contribution evaluation and incentive design. As noted by Recasens et al. \cite{recasens2024frida}, free-rider detection can be framed as an extreme case of the broader contribution-evaluation problem, for which the Shapley value is the game-theoretic gold standard; however, exact computation is combinatorially intractable, and approximations such as Data Shapley and task-specific Shapley valuation are generally designed to \emph{rank} contributors rather than precisely identify zero-contribution participants \cite{ghorbani2019data,jia2019efficient}, limiting their direct applicability to free-rider detection. Log-based auditing schemes detect free riders by inspecting client activity or stake logs rather than model updates directly, while incentive- and reputation-based frameworks allocate rewards or model access proportionally to accumulated trust across rounds. These approaches are complementary to update-based detection but operate on information beyond the raw model updates (e.g., detailed activity logs or reward accounting) and remain vulnerable to strategic clients who behave honestly during an initial trust-building phase before defecting once sufficient reputation has accrued.

\section{Summary and Research Gap}
\label{subsec:research-gap}

Table~\ref{tab:related-work-summary} summarizes the reviewed approaches according to their privacy posture, computational cost, and robustness to adaptive, stealthy free-rider strategies such as that of Zhu et al. \cite{zhu2023advanced}.

\begin{table}[htbp]
\centering
\caption{Comparative summary of existing free-rider mitigation approaches}
\label{tab:related-work-summary}
\begin{tabular}{p{3cm}p{3.6cm}p{2.1cm}p{2.1cm}p{2.4cm}}
\toprule
\textbf{Category} & \textbf{Representative Work} & \textbf{Privacy-Preserving} & \textbf{Lightweight} & \textbf{Robust to Advanced Attacks} \\
\midrule
Foundational \& statistical detection & Fraboni et al. \cite{fraboni2021freerider}; Lin et al. \cite{lin2019freeriders} (STD-DAGMM) & Yes & Moderate & Limited \\
Advanced attack construction & Zhu et al. \cite{zhu2023advanced} & N/A (attack) & N/A & N/A \\
Server-side audit-based & FRAD \cite{wang2024frad}; PASS \cite{wang2023pass} & Yes & Moderate & Limited--Moderate \\
Privacy-attack-based & FRIDA \cite{recasens2024frida} & Partial\footnotemark & No (MIA); Moderate (PIA) & Yes (esp. PIA) \\
Hardware/ cryptographic & TEE, PoW, zero-knowledge proofs & Yes & No & Yes \\
Contribution/ incentive-based & Shapley approximations; reputation systems & Yes & No & Partial \\
\bottomrule
\end{tabular}
\end{table}
\footnotetext{FRIDA requires server-side inspection of individual client updates, in tension with secure aggregation; the authors propose a two-server architecture as a partial mitigation, and show PIA-based detection remains effective under local differential privacy.}

Several patterns emerge from this review. First, purely statistical and feature-based methods \cite{fraboni2021freerider,lin2019freeriders} are lightweight and privacy-preserving but were explicitly shown to be circumventable by attacks constructed to match their target statistics \cite{zhu2023advanced}, and this vulnerability propagates into the audit-based methods (FRAD, PASS) that build on similar anomaly-detection backbones. Second, the privacy-attack-based approach of FRIDA \cite{recasens2024frida} achieves markedly better robustness against exactly this class of advanced, geometry-aware attacks, but its MIA variants require client compliance in training on server-issued canary data (an assumption that a resource-constrained or adversarially adaptive free rider may violate or exploit), and both MIA and PIA variants require the server to inspect individual, non-aggregated client updates—placing FRIDA in direct tension with secure aggregation, one of FL's core privacy-preserving mechanisms. Third, hardware- and cryptography-based methods offer strong guarantees largely independent of the free rider's statistical sophistication, but their cost, hardware dependency, and poor fit for heterogeneous, resource-constrained deployments make them impractical for many real-world FL scenarios. Finally, contribution- and incentive-based methods provide a longitudinal, economically grounded perspective but are either computationally intractable in exact form or vulnerable to strategic, time-varying adversarial behavior.

Across all categories, no reviewed method simultaneously satisfies three properties that are jointly necessary for practical deployment in large-scale, real-world FL systems: (i) \emph{privacy preservation}, in the sense of not requiring the server to inspect raw, individually attributable client updates or client data, and remaining compatible with mechanisms such as secure aggregation; (ii) \emph{lightweight operation}, avoiding the training of auxiliary models (autoencoders, shadow models) or the imposition of additional client-side training and communication burdens that scale poorly with model size or client population; and (iii) \emph{robustness to dynamic, adaptive adversaries}, remaining effective against free-rider strategies explicitly designed—as in Zhu et al. \cite{zhu2023advanced}—to evade previously published defenses, and against client populations and data distributions that shift over the course of training. Methods that are lightweight and privacy-preserving (feature-based statistics) sacrifice robustness to advanced attacks; methods that are robust and reasonably privacy-conscious (FRIDA) sacrifice lightweight operation and full compatibility with secure aggregation; and methods that are robust and lightweight in the security sense (hardware/cryptographic) sacrifice deployability, cost-effectiveness, and applicability to heterogeneous edge and IoT clients.

This unresolved three-way trade-off constitutes the central research gap motivating this thesis. The remainder of this work develops a free-rider detection and mitigation framework explicitly designed to occupy the previously unaddressed region of this design space: relying only on information that is either already available to the server or obtainable without violating standard FL privacy guarantees, avoiding the training of auxiliary detection models or reliance on client-side compliance with server-issued probing data, and evaluated explicitly against adaptive, geometry-aware free-rider strategies across dynamic, non-IID, and large-scale client populations.

\chapter{Proposed Methodology}
\label{sec:proposed-methodology}

\section{Design Objective and Information Boundary}

The proposed method is a server-orchestrated challenge--response mechanism for free-rider mitigation. The server occasionally sends a selected client a randomized model in place of the current global model. This secret model is called a \emph{trap}. An honest client applies the prescribed local optimizer to whichever model it receives, whereas a free rider fabricates a response from its locally visible model history. Because the client is never told whether a received model is a trap, the intervention changes the attacker's input without changing the client API.

The detector trusts neither local accuracy nor loss, reported training time, hardware counters, number of steps, nor any other client-supplied claim. Its decision inputs are derived by the server from the returned model: whether the update is effectively zero, the norm of the flattened update, and the relative distribution of update magnitude across parameter tensors. Ground-truth free-rider labels and the configured attack name are used only to evaluate the simulator; they are not available to the detection rule. Cosine similarity and local metrics are logged for analysis but do not affect flags, penalties, state transitions, anchor selection, or removal.

This boundary makes the method lighter than an auxiliary autoencoder or a client-side proof-of-work mechanism. It does not, however, constitute a cryptographic proof of training. It also assumes that the server can inspect an individual returned model during an audit; ordinary secure aggregation would require a separate auditable channel or a privacy-preserving computation of the required statistics.

\section{System and Threat Model}

Let $N$ be the registered clients and $T$ the communication rounds. Client $i$ receives $w^{\mathrm{send}}_{t,i}$ in round $t$ and returns $w^{\mathrm{ret}}_{t,i}$. The server derives
\begin{equation}
\Delta_{t,i}=w^{\mathrm{ret}}_{t,i}-w^{\mathrm{send}}_{t,i},
\qquad r_{t,i}=\lVert\Delta_{t,i}\rVert_2.
\end{equation}
In an ordinary round $w^{\mathrm{send}}_{t,i}=w^g_t$; a trapped client instead receives $w^{\mathrm{trap}}$. Trap-derived updates are never aggregated because they do not share the global starting point.

The implementation evaluates four attack behaviors through the same server rule:
\begin{itemize}
    \item \textbf{FR1---replay:} return the last model observed by that client; on first participation, return the current received model.
    \item \textbf{FR2---bounded random update:} add independent uniform noise in a fixed small range to each floating-point parameter of the received model.
    \item \textbf{FR3---historical average:} average the latest five models received by that client, including any unlabelled traps.
    \item \textbf{FR4---trajectory prediction:} estimate an expected update from the latest five differences between consecutively received models and add Gaussian noise equal to 10\% of the estimated norm. Before history exists, use a small random cold-start update.
\end{itemize}
The attacker sees its identity, round number, received model, public configuration, its own received-model history, and local randomness. It receives no trap indicator. Collusion, defense-aware adaptation, on--off participation, and partially honest local training are outside the evaluated scope.

\section{Robust Two-Signal Detector}

\subsection{Two-sided magnitude evidence}
For a reference set of update norms $R$, the server computes
\begin{align}
m_r &= \operatorname{median}(R),\\
d_r &= \operatorname{median}_{r\in R}|r-m_r|,\\
z^{\mathrm{norm}}_{t,i} &= \frac{r_{t,i}-m_r}{\max(d_r,0.05)}.
\end{align}
The magnitude signal fires when $|z^{\mathrm{norm}}_{t,i}|>3$. The absolute value is essential: replay may create an excessively large response, whereas an adaptive fabricated update may be abnormally small. The MAD floor prevents a nearly identical reference population from turning negligible numerical differences into extreme standardized scores.

\subsection{Scale-free layer profile}
For each floating parameter tensor $k$, the server forms
\begin{equation}
p_{t,i,k}=\frac{\lVert\Delta_{t,i,k}\rVert_2}
{\sum_j\lVert\Delta_{t,i,j}\rVert_2}.
\end{equation}
The vector $p_{t,i}$ records how update energy is distributed across layers independently of total scale. A coordinate-wise median reference profile is constructed from the applicable comparison population. The Euclidean distance from this center is standardized by median and MAD, with a floor of $0.01$; the profile signal fires above 3.

A single magnitude or profile anomaly is sufficient to quarantine an unflagged client for confirmation, but it contributes no magnitude penalty. Only a response that fails both tests on the same probe creates \emph{joint gradient evidence}. This conjunction was introduced because persistent honest clients can occupy the tail of either one-dimensional statistic.

\subsection{Exact send-back evidence}
If $r_{t,i}<10^{-6}$, the response is treated as an exact or near-exact send-back. It immediately adds five permanent points and sends the client directly to the suspicious state. This check operates in every phase, including warm-up. It is evaluated separately from nonzero magnitude/profile evidence and is never erased by rehabilitation.

\section{Client States and Evidence Accounting}

Each active client occupies one of three operational states:
\begin{description}
    \item[Unflagged ($U$)] no unresolved anomaly; the update is normally eligible for aggregation unless trapped in the current round.
    \item[Candidate ($C$)] at least one nonzero anomaly has occurred and the client awaits fresh-trap confirmation. Candidate updates are quarantined.
    \item[Suspicious ($S$)] repeated joint evidence or an exact send-back has confirmed elevated suspicion. The client is probed frequently.
\end{description}

The fixed removal threshold is 15 points. A joint anomaly contributes two points in initial coverage, one point in later surveillance, and three points in candidate/suspicion probing. A magnitude-only or profile-only anomaly contributes zero points. Nonzero evidence cannot remove a client until at least ten relevant episode probes have been observed. Three exact send-backs reach 15 points independently of that gate.

For $C$ and $S$ clients, the server maintains episode probe and joint-failure counts. After at least ten probes, a client is rehabilitated when its joint-failure rate is at most 10\%. Rehabilitation returns it to $U$, clears the current magnitude episode and resets episode counters, while retaining exact-send-back points. If rehabilitation and nonzero removal become eligible simultaneously, rehabilitation is applied first.

\begin{table}[htbp]
\centering
\caption{Principal protocol parameters}
\label{tab:v8-parameters}
\begin{tabular}{lll}
\toprule
Parameter & Meaning & Value \\
\midrule
$n_w$ & warm-up rounds & 10 \\
$\rho$ & nominal trap group fraction & 0.10 \\
$k_r,k_p$ & norm/profile thresholds & 3.0, 3.0 \\
$\epsilon_r,\epsilon_p$ & norm/profile MAD floors & 0.05, 0.01 \\
$\delta_0$ & exact send-back penalty & 5 \\
$\delta_{\mathrm{cov}}$ & initial joint penalty & 2 \\
$\delta_Y$ & surveillance joint penalty & 1 \\
$\delta_S$ & suspicion joint penalty & 3 \\
$P_{\max}$ & removal threshold & 15 \\
$q_{\min}$ & rehabilitation probe minimum & 10 \\
$\eta_{\max}$ & rehabilitation failure rate & 0.10 \\
\bottomrule
\end{tabular}
\end{table}

\section{Probing Schedule}

\subsection{Warm-up}
Rounds 1--10 use standard FedAvg. The exact send-back test remains active, but nonzero magnitude/profile scoring is disabled so early optimization transients do not create candidates.

\subsection{Initial double coverage sweep}
After warm-up, clients are shuffled into groups of approximately 10\%. The server generates two distinct traps, A and B, from a global base. Each trap is frozen for one complete pass. Two groups are probed per round: a randomized first pass gives every client trap A once, and an independently randomized second pass gives every client trap B once, avoiding repeated group pairs where possible. Thus every client receives two secret probes during ten coverage rounds.

Responses are buffered until a full pass is complete. Statistics are computed separately for A and B, ensuring that responses to different traps do not define one another's baseline. Exact send-backs add five points each; each joint norm/profile failure adds two; a one-signal anomaly creates candidate evidence but no magnitude points. Ambiguous clients enter $C$, while sufficiently repeated evidence enters $S$. The initial double sweep is performed once.

\subsection{Post-coverage surveillance}
All currently unflagged clients are shuffled into exactly ten $Y$ groups. One group is audited per round using the cycle's frozen surveillance trap and excluded from that round's aggregation. After all groups are visited, a new shuffled cycle begins. Any anomaly sends a client to $C$ and permanently disqualifies it from anchor service; only a joint anomaly adds one point.

\subsection{Candidate and suspicion probing}
When $S$ is nonempty, the number of suspicion groups is
\begin{equation}
G_S=\max\left(1,\left\lfloor\frac{|S|}{10}\right\rfloor\right).
\end{equation}
Ten is a target, not a maximum group size. Candidates are distributed across the same groups, or one confirmation group is formed when only candidates exist. The implementation reshuffles suspicious and candidate membership across groups each round. Each group receives a distinct fresh trap.

Each group receives up to three active anchors, subject to a system-wide maximum of ten. Suspect and candidate responses are compared only with anchors that received the same trap. Anchor-relative medians and MADs define both decisions. A joint failure adds three points; a normal candidate response clears $C$ immediately. Two joint failures confirm a candidate as suspicious by default.

\section{Rotating Anchor Roster}

The v8 change replaces fixed panels with a rotation queue. Eligible anchors are active, unflagged clients whose lifetime anomaly count is zero. An anchor is not reused until the eligible queue has been traversed where population size permits. Clients in the current $Y$ group cannot simultaneously anchor. Off-duty roster members receive the ordinary global model and contribute normally; only active anchors receive a trap and are excluded.

Rotation reduces repeated comparison against one unlucky panel, and $Y$ auditing can remove an anomalous client from future eligibility. Nevertheless, eligibility is absence of evidence, not proof of honesty. Chapter~\ref{sec:results-discussion} shows that adaptive attackers can evade surveillance, enter the roster, and contaminate panels.

\section{Aggregation, Removal, and Checkpointing}

Every trap recipient---suspect, candidate, anchor, or $Y$ client---is excluded from aggregation in that round. Candidates remain quarantined through confirmation. Permanent removal occurs when exact-send-back evidence reaches 15, or when total retained points reach 15 after the ten-probe nonzero-evidence gate and rehabilitation does not apply.

The simulator saves a checkpoint after every completed round, retains the newest five while a run is incomplete, and deletes checkpoints after successful completion. Checkpoints include model and attacker-visible history, client states, penalties, episode counters, buffered coverage evidence, and anchor rotation. Logs record model utility, round time, memory, trap assignments, state transitions, removals, rehabilitation, and per-client detection statistics.

\section{Computational Characteristics and Limitations}

For $d$ model parameters, deriving the norm and profile costs $O(d)$ per inspected update. Median/MAD construction over $m$ observations costs $O(m\log m)$ by sorting and $O(m)$ temporary scalar storage. Persistent state is $O(N)$. No detector model is trained and the client performs only its ordinary optimizer on the model received. The main utility cost is lost aggregation capacity from trap recipients and candidates.

The method observes statistical consequences rather than proving training. Median/MAD has a theoretical 50\% scalar breakdown point, but practical separation can fail earlier when honest non-IID responses are broad or multimodal. Anchor rotation cannot help if undetected attackers qualify as anchors. Per-client inspection conflicts with ordinary secure aggregation. Evaluation is limited to MNIST and CIFAR-10, two primary seeds, two attacker shares, full participation, and always-on non-colluding attacks; dropout, collusion, on--off behavior, selfish partial training, and defense-aware adaptation remain future work.

\section{Summary}

The proposed method combines secret two-trap coverage, two-sided norm and layer-profile evidence, candidate confirmation, frequent suspicion probing, rotating anchors, continued surveillance, cumulative penalties, and rehabilitation. Its rules are attack-independent and require no trusted client-reported metric. The next chapter evaluates both its strong IID performance and the limits revealed by heterogeneous data and adaptive FR4 behavior.




\begin{comment}

\chapter{Proposed Methodology} \label{chapter:proposed}
The Proposed Methodology chapter is the core of your thesis, where the approach to solving the research problem is presented in a detailed and structured manner. This chapter aims to provide a clear, step-by-step explanation of how you plan to address the problem outlined in the introduction. Each component of the methodology should be explained in a way that allows readers, even those less familiar with the specific technical details of the project, to understand both the structure and the rationale behind your choices. Clarity is essential, and visual aids should be used to enhance understanding where necessary.

\section{Overview of the Methodology}

To begin, provide a high-level overview of the methodology. This section introduces the key components and offers a broad structure of the overall approach. It is essential to explain why each component was chosen, highlighting its advantages and explaining why it was preferred over other alternatives. The reader should be able to see how these individual components fit together to form a cohesive system aimed at solving the research problem. A comprehensive figure, such as a flowchart or architecture diagram, should be included to visually represent the methodology, serving as a roadmap that helps the reader visualize how the components interact with one another and connect to the broader objectives of the study.


\section{Chronological Breakdown of Components}

Next, break down the methodology into its individual components, presented in the order in which they are executed. This section should offer a logical flow from one component to the next, demonstrating how each part builds upon the previous one to create an integrated system. For each component, several sub-sections will be required.

Begin by clearly describing the role of each component. Explain the specific problem it solves and how it contributes to the overall system. Follow this by discussing the method or algorithm chosen for this component. For example, whether it's a machine learning algorithm, a database structure, or a communication protocol, provide enough technical detail for readers to understand how it works. Reference any theoretical foundations or literature that support your choice, and if a standard algorithm is used, explain why this particular approach was selected over others.

Additionally, specify the type of data required as input for the component and what it produces as output. Where applicable, visual aids such as tables or diagrams should be included to clarify the flow of data. You should also explain the key principles or logic underlying the functionality of the component. For technical algorithms, this might involve discussing higher-level concepts such as how a neural network propagates data forward during training, but without delving into code-level specifics.

Lastly, justify your design decisions. Detail why you chose this particular method or approach. Were there alternatives, and if so, why were they not selected? This section should reflect critical thinking, showing that choices were made deliberately, based on a reasoned evaluation of trade-offs. It's important to acknowledge that sometimes an approach is selected because it is the best available, even if it isn't perfect.

\section{Integration and Interconnections}

In this section, discuss how the individual components interact with each other. The reader should understand how data flows between components, and how these interactions contribute to achieving the overall research objective. This is crucial for demonstrating that the methodology is not just a collection of independent modules but a cohesive system where each component plays an interconnected role. For example, if your system involves feedback loops, dependencies, or iteration processes—such as updating a machine learning model during training—these interactions must be clearly described.

Refer back to the architectural diagram introduced at the beginning of the chapter as you discuss these interactions. This ensures that readers can relate each described component to the larger structure, reinforcing their understanding of the entire methodology.

\section{Clarity and Comprehensiveness in Visual Aids}

Effective use of visual aids—such as diagrams, flowcharts, and graphs—is essential for making complex ideas more understandable. Each visual should add clear value to the explanation and should be meticulously labeled and accompanied by detailed captions that can stand alone. For instance, if your method involves multiple steps in data preprocessing, a flowchart can simplify what might otherwise seem like a complex series of operations.

Aim for clarity in every visual. Avoid clutter, and ensure that each element is well-explained in the text. Good visuals should simplify technical concepts and make them accessible to a broad audience, even for particularly intricate methodologies.

\section{Conclude with a Summary of the Methodology}

Conclude this chapter with a summary of the overall methodology, reiterating how the individual components work together to achieve the research objectives. This recap should give readers a clear understanding of the entire system before moving on to the subsequent chapters on Results and Discussion. The summary serves as a final reminder of the methodology's coherence and structure, providing closure to the technical explanation of your approach.

By following this structure, the Proposed Methodology chapter will present your approach in a clear, well-organized manner, making it accessible and engaging to your readers, while maintaining technical rigor.
\end{comment}

\chapter{Results and Discussion}
\label{sec:results-discussion}

\section{Experimental Questions}

The evaluation addresses four questions. First, can SWT-CP detect replay, random-update, historical-average, and adaptive free riders without knowing which attack is present? Second, how sensitive is detection to the random seed and attacker proportion? Third, does Dirichlet non-IID partitioning increase false positives or create false negatives? Fourth, what model-utility and execution costs accompany the probing schedule?

\section{Experimental Setup}

All reported v8 experiments use MNIST, 100 clients, full participation, 100 communication rounds, three local epochs, batch size 32, SGD with learning rate 0.01 and momentum 0.9, and a convolutional classifier. IID partitions distribute examples uniformly. Non-IID partitions use a Dirichlet label allocation with concentration $\alpha=0.5$. The matrix contains attacker shares of 30\% and 40\%, seeds 42 and 43, and FR1--FR4, yielding 32 completed runs. No result from v4--v7 is pooled with v8.

Precision, recall, and F1 treat permanent removal as the positive decision. Thus $TP$ is a removed free rider, $FP$ a removed honest client, and $FN$ an active free rider at round 100. Final candidates are not counted as false removals. This distinction is important because candidates created in round 100 have no later confirmation round. Global accuracy is evaluated on the MNIST test set. Resource measurements are process-level measurements from the simulator; resumed-run wall-clock totals may describe only the resumed segment, so average round time is used for cross-run comparison.

\begin{table}[htbp]
\centering
\caption{Aggregate v8 results across FR1--FR4 for each setting}
\label{tab:aggregate-results}
\resizebox{\textwidth}{!}{%
\begin{tabular}{llrrrrrrr}
\toprule
Partition & Attackers/seed & TP & FP & FN & Precision & Recall & F1 & Mean accuracy \\
\midrule
IID & 30\%/42 & 120 & 8 & 0 & 93.75\% & 100\% & 96.77\% & 98.95\% \\
IID & 30\%/43 & 120 & 1 & 0 & 99.17\% & 100\% & 99.59\% & 98.92\% \\
IID & 40\%/42 & 160 & 6 & 0 & 96.39\% & 100\% & 98.16\% & 98.86\% \\
IID & 40\%/43 & 160 & 0 & 0 & 100\% & 100\% & 100\% & 98.89\% \\
Non-IID & 30\%/42 & 120 & 10 & 0 & 92.31\% & 100\% & 96.00\% & 98.88\% \\
Non-IID & 30\%/43 & 120 & 9 & 0 & 93.02\% & 100\% & 96.39\% & 98.88\% \\
Non-IID & 40\%/42 & 120 & 7 & 40 & 94.49\% & 75\% & 83.62\% & 98.84\% \\
Non-IID & 40\%/43 & 120 & 8 & 40 & 93.75\% & 75\% & 83.33\% & 98.88\% \\
\bottomrule
\end{tabular}}
\end{table}

\section{IID Results}

All 560 free riders in the four IID configurations were removed. Seed 43 was particularly strong: one honest client was removed across the 30\% experiments and none across the 40\% experiments. Seed 42 produced eight and six honest removals respectively. The gap shows that a single seed would give a misleading estimate of precision even under IID partitioning. Across all IID runs, pooled removal precision was $560/(560+15)=97.39\%$, recall was 100\%, and F1 was 98.68\%.

FR1--FR4 all remained detectable under IID data. Final accuracy varied by only 0.14 percentage points across the aggregate settings. The defense therefore preserved utility in these experiments despite excluding trapped and candidate updates. However, the seed-42 false removals demonstrate that joint evidence and rotation reduce rather than eliminate honest tails.

\section{Non-IID Results}

Table~\ref{tab:noniid-detail} presents each non-IID attack as ``free riders removed / honest clients removed.'' FR1--FR3 achieved complete recall in every setting. FR3 was the least precise, producing three or four false removals per run. FR4 exhibited a qualitatively different transition: it was fully detected at 30\% but completely missed at 40\% under both seeds.

\begin{table}[htbp]
\centering
\caption{Detailed non-IID results (free riders removed / honest clients removed)}
\label{tab:noniid-detail}
\begin{tabular}{lrrrr}
\toprule
Attack & 30\%/seed 42 & 30\%/seed 43 & 40\%/seed 42 & 40\%/seed 43 \\
\midrule
FR1 & 30/3 & 30/1 & 40/1 & 40/3 \\
FR2 & 30/2 & 30/3 & 40/1 & 40/2 \\
FR3 & 30/3 & 30/3 & 40/4 & 40/3 \\
FR4 & 30/2 & 30/2 & 0/1 & 0/0 \\
\bottomrule
\end{tabular}
\end{table}

At 30\%, all FR4 clients entered the candidate/suspicion pipeline and were removed around rounds 30--32. Joint evidence appeared in 76.01\% of trapped FR4 observations for seed 42 and 81.28\% for seed 43. At 40\%, the corresponding rates fell to 0\% and 2.86\%; no attacker accumulated enough repeated evidence. Seed 42 never flagged any of its 40 FR4 clients. Seed 43 flagged 18 at least once, but none received adequate confirmation. Every one of the 80 attackers across the two 40\% FR4 runs remained active at round 100.

This discontinuity is visible during initial coverage. Median FR4 norm z-scores were $-2.65$ and $-3.10$ at 30\%, but only $-1.55$ and $-1.81$ at 40\%. Increasing contamination changed the median/MAD and layer-profile references enough to make the malicious cluster appear less exceptional, even though its raw update norms remained much smaller. In the seed-42 40\% run, the median trapped norm was 0.0365 for FR4 and 0.4824 for honest clients, a factor of approximately 13. The detector failed not because responses were physically similar, but because the reference and joint-decision mechanism did not convert that separation into evidence under the broad non-IID mixture.

\section{Anchor Contamination Case Study}

The 40\% FR4 failure also invalidated the operational interpretation of an eligible anchor. Because no FR4 client had prior evidence, attackers entered the never-flagged roster. In the seed-42 run, honest client 85 became suspicious and was compared with 33 anchor assignments over its episode; 10 were free riders. At rounds 66, 68, and 73, two of its three anchors were FR4 clients. Those exact three malicious-majority panels produced joint norm/profile failures and added three points each. Combined with six earlier points, client 85 reached 15 and was removed at round 73.

This trace establishes a causal failure rather than a general suspicion: rotating an unverified roster distributes contamination but does not remove it. Three anchors tolerate one bad reference, not two. A client that has never been detected is not necessarily a positively verified honest client.

\section{False Positives and Seed-Specific Clients}

Non-IID false removals were concentrated rather than uniformly random. Under seed 42, client 65 was removed in six non-IID attack/configuration combinations and client 74 in five. Under seed 43, client 81 was removed seven times and client 29 five times. Repetition across FR1--FR4 indicates partition-specific honest update geometry rather than an attack-specific response. The detector repeatedly interprets certain legitimate label distributions as joint norm/profile anomalies.

At round 100 the non-IID settings contained between 17 and 34 candidates and one to four suspects when aggregated across the four attacks. Most final candidates entered in the last round and are boundary artifacts rather than confirmed removals. Final suspects are more important: they show that removal precision alone understates temporary quarantine and future false-removal risk.

\section{Model Utility}

Mean final accuracy remained between 98.84\% and 98.95\% in every aggregate setting. The individual-run range was 98.82--98.95\%. Notably, the non-IID 40\% FR4 runs still achieved 98.84\% and 98.88\% despite detecting no attackers. FR4 returns low-magnitude fabricated updates; with 60 honest clients remaining, these responses mainly dilute learning rather than immediately destroy classification accuracy. Consequently, high test accuracy is not evidence of successful free-rider detection.

\section{Runtime and Resource Use}

\begin{table}[htbp]
\centering
\caption{Execution profile aggregated across four attacks}
\label{tab:resources}
\begin{tabular}{llrrr}
\toprule
Partition & Attackers/seed & Mean s/round & Approx. batch hours & Peak RAM (MB) \\
\midrule
IID & 30\%/42 & 50.8 & 5.65 & 1074 \\
IID & 30\%/43 & 43.4 & 4.82 & 1150 \\
IID & 40\%/42 & 59.6 & 6.62 & 3172 \\
IID & 40\%/43 & 54.3 & 6.04 & 3237 \\
Non-IID & 30\%/42 & 46.9 & 5.21 & 887 \\
Non-IID & 30\%/43 & 46.0 & 5.11 & 891 \\
Non-IID & 40\%/42 & 42.4 & 4.72 & 1079 \\
Non-IID & 40\%/43 & 44.8 & 4.98 & 1022 \\
\bottomrule
\end{tabular}
\end{table}

The IID 40\% runs recorded CUDA memory, while the other collected batches executed on CPU. GPU execution did not yield proportional speedup because clients are simulated sequentially and much of the workload and orchestration remains serial. FR4 40\% non-IID was relatively costly within its batch because all attackers remained active: the seed-42 run averaged 89 aggregated clients per round and ended with 99 active clients, compared with roughly 55--56 aggregated and 56--59 active clients after successful FR1--FR3 removal. Peak recorded system-memory use approached 95\% in some CPU runs, so host memory was a practical execution constraint.

\section{Threats to Validity}

Only two seeds were evaluated, and their false-positive difference is large. The experiments use one dataset, one model family, one Dirichlet concentration, full participation, and always-on independent attackers. Reported percentages are therefore measurements of this simulation, not confidence intervals for deployment. No external defense was rerun under an identical implementation, so the results do not establish superiority over FRIDA, FRAD, PASS, or STD-DAGMM. Accuracy remains high even in a detection failure, and should not be treated as a proxy for contribution authenticity. Finally, individual-update inspection is incompatible with standard secure aggregation without additional protocol support.

\section{Overall Interpretation}

The central empirical conclusion is conditional. SWT-CP v8 is highly effective against FR1--FR3 across the tested IID and non-IID settings, and against FR4 under IID or 30\% non-IID contamination. It is not robust to FR4 at 40\% non-IID contamination. The repeated two-seed failure indicates a methodological boundary rather than a corrupted run. Future revisions should require positive evidence before anchor admission and exploit paired within-client responses to the two coverage traps, reducing dependence on comparisons among naturally heterogeneous clients.

\chapter{Conclusion}

This thesis investigated whether a server can mitigate free riding using only returned model updates, without trusting client-reported accuracy, loss, training time, or resource use and without training an auxiliary detector. SWT-CP v8 addresses this objective through secret trap models, two-sided norm and layer-profile evidence, a candidate/suspicion state machine, repeated probing, rotating anchors, cumulative penalties, and rehabilitation.

The IID evaluation is encouraging: every tested free rider was removed across attacker shares of 30\% and 40\% and seeds 42 and 43. Seed 43 produced only one honest removal at 30\% and none at 40\%. FR1--FR3 also retained complete recall under non-IID partitioning. These results support the use of unpredictable server interventions and longitudinal evidence for lightweight detection of replay, random-update, and historical-average behavior.

The evaluation also revealed two limits that are as important as the successes. First, non-IID false positives concentrate on particular honest clients and recur across attacks, showing that joint statistical evidence is still sensitive to legitimate client heterogeneity. Second, FR4 exhibits a reproducible contamination transition: both non-IID 30\% runs removed all attackers, whereas both 40\% runs removed none. Undetected FR4 clients entered the anchor roster, and malicious-majority panels directly contributed to an honest removal. Rotation is therefore not equivalent to trust establishment.

The immediate research priority is to replace negative anchor eligibility (``never flagged'') with positive, repeated validation, and to use differential responses from the same client across independent traps so that each client partly serves as its own non-IID control. A third seed and intermediate 35\% FR4 experiment should confirm the transition. Broader validation should include CIFAR-10, additional Dirichlet concentrations, client dropout, on--off and partial-training free riders, collusion, defense-aware adaptation, and comparisons with implemented baselines. A secure protocol is also needed if per-client auditing is to coexist with secure aggregation.

SWT-CP should thus be understood as a computationally simple and empirically strong detector within a defined operating region, not as proof that a client trained. Its main contribution is both constructive and diagnostic: it demonstrates where secret statistical traps work, and identifies the precise combination of heterogeneity, adaptive fabrication, and contaminated references under which they fail.


\begin{comment}
    
\section{Datasets and Experimental Setup}

The chapter begins with a clear description of the datasets used in your research. It's crucial to provide detailed information about the features of the data, including whether it is text-based, numerical, image-based, or of another type. You should also explain the size of the dataset, its sample distribution, and any inherent biases or imbalances that may affect the analysis. For instance, if your dataset has a significant class imbalance or contains noisy or missing data, describe the strategies employed to mitigate these issues, such as data augmentation, imputation methods, or rebalancing techniques. Including sample data or visualizations can further clarify the nature of the dataset, offering a glimpse into the raw materials of your research. This could take the form of snippets of text, feature distributions, or representative images, depending on the context.

Once the dataset has been outlined, provide a thorough description of the experimental setup. Detail the computational environment, including the hardware, software libraries, and any cloud services utilized. Key hyperparameters such as learning rate, batch size, and the number of epochs should be listed, and if possible, their selection should be justified either theoretically or empirically. This justification could involve citing similar studies or explaining any tuning processes you conducted. Offering these insights helps the reader understand why specific choices were made and how they might influence the results.

\section{Presenting the Results}

The results themselves are presented next, starting with a brief summary of what types of outcomes you will report and how they relate to your research objectives. To ensure clarity and precision, results should be presented systematically. For quantitative findings, such as performance metrics (e.g., accuracy, precision, recall, F1 score), tables are an effective medium. Each table should be clearly labeled and referenced in the text, with concise descriptions of the data they present.

Graphical representations such as bar charts, line graphs, and scatter plots should be used to illustrate trends or comparisons, such as how different models perform over time or how various configurations influence results. When appropriate, confusion matrices or ROC curves can be employed to visualize classification performance in more detail. However, at this stage, your focus should be on reporting the data, not interpreting it. If your research involves multiple phases or experiments, presenting each experiment's results separately will make the data easier to follow, especially when the phases involve distinct datasets or algorithms.

\section{Interpreting the Results}

Once the raw data has been presented, the focus shifts to interpreting the findings. This involves analyzing how the results align with or diverge from the hypotheses or research questions posed earlier. If certain findings contradict expectations, possible explanations should be provided, supported by relevant theory or prior research. For instance, unexpected results might be due to limitations in the dataset, overfitting, or other factors.

Strengthen your interpretations by linking the results back to theoretical frameworks or previous studies. If a particular algorithm outperforms others, consider citing earlier work that supports this outcome, explaining why this might be the case. For example, you might discuss how the algorithm generalizes better due to a superior feature extraction process or faster convergence rates. On the other hand, if certain results are disappointing, explore the potential causes and consider whether similar patterns have been observed in related research.

Comparisons with similar studies are important in this section. Highlight similarities and differences between your results and those reported by others, offering possible reasons for these divergences. For instance, you might say, "While Author A's method achieved an accuracy of 85% on a similar dataset, our approach improved this to 90% by employing an additional feature extraction step."

\section{Challenges and Limitations}

Reflecting on the challenges and limitations encountered during your research is an essential part of the discussion. Consider how these challenges may have impacted your results. Were there limitations in the dataset, such as a lack of diversity or imbalance? Did methodological or computational constraints affect the experiments? This reflection not only adds depth to your analysis but also paves the way for future research by acknowledging areas where further improvements could be made. Additionally, recognizing these limitations allows readers to contextualize the results within the scope of your study.

\section{Implications and Significance of Findings}

The broader implications of your findings need to be explored. What new knowledge or insights does your research contribute to the field? Consider how your results may influence future work, whether by suggesting new directions for research, optimizing existing methods, or offering real-world applications. For example, if your findings improve neural network performance in a specific domain, could this optimization be applied to industry settings, such as more efficient data processing or improved cybersecurity measures? By positioning your findings within the wider landscape of the field, you underscore the relevance and potential impact of your work.

\section{Decision to Split or Combine Results and Discussion}

Depending on the nature of your results, you may choose to either combine or split the Results and Discussion sections. If your research has produced a large volume of complex data, it may be beneficial to split these sections. In such cases, the Results chapter would focus solely on presenting the raw data, while the Discussion chapter would handle the analysis and interpretation. This approach can also be helpful if your research involved multiple distinct phases or both quantitative and qualitative data, allowing you to deal with each type or phase independently.

On the other hand, if your results are straightforward or closely tied to their interpretation, combining the two sections can lead to a more concise and cohesive presentation. When results and interpretations are tightly intertwined—where explaining trends or patterns requires immediate analysis—a combined chapter helps maintain a logical flow, making it easier for the reader to follow the narrative.

\section{Conclusion of the Chapter}

Conclude this chapter by summarizing the key takeaways from your results and their significance. This final section should tie everything together, leading naturally into the Conclusion chapter, where you will reflect on the broader contributions of your research and suggest potential directions for future work. The summary serves as a moment of synthesis, reinforcing the most important insights gained from your findings and setting the stage for the overarching conclusions of your thesis.

This structure ensures that the Results and Discussion chapter not only presents data but also thoughtfully interprets it, giving readers a comprehensive understanding of how the research findings relate to the broader field. By organizing the chapter in this manner, you maintain clarity and coherence, making the technical details accessible and engaging.


\chapter{Conclusion}
The Conclusion chapter is a critical reflection on the research conducted, encapsulating the key findings, contributions, and broader implications of the work. This final chapter serves to tie together the insights gained from your study, while also addressing its limitations and proposing potential directions for future research. The tone should be reflective, confident, and forward-looking, emphasizing the significance of the work without overstating its claims.

\section{Restating Research Objectives and Questions}

To begin, it is important to restate the main research objectives and questions that framed your thesis. This serves as a reminder of the foundational focus of the study, bringing the reader back to the core issues you set out to address. You should rephrase these objectives and questions succinctly, demonstrating how they guided the entire research process and structured the work you have accomplished. By revisiting these questions, you provide a clear context for evaluating the outcomes and contributions of your research.

\section{Summary of Key Findings}

Following the restatement of your research objectives, you should offer a concise summary of the most significant findings from your study. Focus on highlighting the key results that directly address your research objectives, ensuring that the discussion remains tightly connected to the goals outlined at the beginning of the thesis. It is essential to present the insights that stand out the most, emphasizing those that contribute the most meaningfully to the research questions posed. This summary should distill the essence of your findings, providing a clear and accessible synthesis of the work.

\section{Discussion of Implications}

Next, reflect on the broader implications of your findings. Consider how your results contribute to the existing body of knowledge within your field and how they answer the research questions you initially posed. This discussion should explore any theoretical, methodological, or empirical contributions made by your study. For example, if your findings offer new insights into a particular phenomenon or propose a novel methodological approach, these contributions should be clearly articulated. Moreover, it is important to relate your findings to the wider academic discourse, connecting them to the works discussed in the Related Works chapter. By positioning your research within the broader field, you highlight its relevance and the value it brings to ongoing scholarly conversations.

\section{Contributions of the Research}

You should provide a clear and structured account of the contributions your research has made. These contributions may take various forms: theoretical insights that challenge or extend existing frameworks, practical applications that could influence real-world scenarios, methodological advancements that provide new tools for future research, or empirical findings that add to the data-driven understanding of your field. Each type of contribution should be acknowledged and clearly articulated, giving readers a sense of the unique value your work brings to the academic and professional landscape.

\section{Acknowledging Limitations}

No research is without limitations, and it is important to be transparent about the constraints that may have influenced your results. Discuss the specific limitations of your study, such as sample size, scope, data quality, or technical challenges that arose during the research process. Acknowledging these limitations not only demonstrates intellectual honesty but also helps contextualize your findings, showing how they fit within the constraints of the study. Moreover, by openly discussing the limitations, you provide a foundation for future researchers.

\section{Suggestions for Future Research}

Based on your findings and the limitations identified, you should offer clear and thoughtful recommendations for future research. This section provides an opportunity to identify specific areas where further investigation is needed, suggesting how future studies might address unanswered questions or overcome the limitations you encountered. Where relevant, you may also propose new research directions or methodologies that could be applied to similar problems, offering a roadmap for how the field can advance in light of your work. By pointing to areas that require further exploration, you not only acknowledge the ongoing nature of academic inquiry but also position your research as a stepping stone for future developments.

\section{Final Reflections}

As you approach the end of the chapter, a final reflection on your research journey is appropriate. This could be a brief statement that encapsulates the significance of your study and its place within the broader research community. You may wish to share personal insights gained from conducting the research, including any unexpected outcomes or challenges that led to deeper learning. This reflection serves to humanize the research process, providing a glimpse into the intellectual and practical growth experienced throughout the thesis. It is also an opportunity to reflect on the broader impact of the study, reinforcing the relevance and potential contributions of your work.

\section{Concluding Remarks}

To close the chapter, end with a strong and definitive statement that encapsulates the contribution of your research and its potential for future impact. This final statement should provide a sense of closure, leaving the reader with a clear understanding of the significance of your work. Whether the research has practical implications, theoretical advancements, or methodological contributions, the concluding remarks should confidently assert the value of the study and its place in the field, while also gesturing toward the future possibilities it opens up.

This approach ensures that the Conclusion chapter not only summarizes the research but also thoughtfully reflects on its significance and future potential, leaving the reader with a clear and compelling understanding of the work's broader impact.

\end{comment}

%-------------------------- BACK MATTER --------------------------%

% Print the references used throughout the document, mandatory
\printbib

\end{document}
